% One-page resume for SWE internship applications. Source of truth for facts is
% the master resume; keep the two in sync. Build: tectonic resume/resume.tex
% then move resume/resume.pdf to the repo root (see README).
\documentclass[10pt,letterpaper]{article}

\usepackage[margin=0.5in]{geometry}
\usepackage{enumitem}
\usepackage{titlesec}
\usepackage[hidelinks,pdftitle={Ryan Ye - Resume},pdfauthor={Ryan Ye},pdfsubject={Computer Science, Cornell University}]{hyperref}
\usepackage{xcolor}
\usepackage[T1]{fontenc}
\usepackage{lmodern}

\pagestyle{empty}
\setlength{\parindent}{0pt}

\titleformat{\section}{\normalsize\bfseries\scshape}{}{0pt}{}[\vspace{-2pt}\titlerule]
\titlespacing{\section}{0pt}{6pt}{3pt}

\newcommand{\entry}[4]{%
  \noindent\begin{tabular*}{\textwidth}{@{}l@{\extracolsep{\fill}}r@{}}
    \textbf{#1} & #2 \\
    \textit{#3} & \textit{#4} \\
  \end{tabular*}\vspace{-1pt}}

\newcommand{\entryshort}[2]{%
  \noindent\begin{tabular*}{\textwidth}{@{}l@{\extracolsep{\fill}}r@{}}
    \textbf{#1} & #2 \\
  \end{tabular*}\vspace{-1pt}}

\setlist[itemize]{leftmargin=1.2em, itemsep=0pt, topsep=1pt, parsep=0pt, partopsep=0pt}

\begin{document}

\begin{center}
  {\LARGE\bfseries Ryan Ye}\\[2pt]
  \href{mailto:rmy43@cornell.edu}{rmy43@cornell.edu} $\cdot$
  \href{https://ryanmye.github.io}{ryanmye.github.io} $\cdot$
  \href{https://www.linkedin.com/in/rmy43}{linkedin.com/in/rmy43} $\cdot$
  \href{https://github.com/ryanmye}{github.com/ryanmye}
\end{center}

\section{Education}
\entry{Cornell University}{Ithaca, NY}{B.S. in Computer Science, College of Engineering, GPA 4.02}{Expected May 2028}
\begin{itemize}
  \item \textbf{Coursework:} Deep Learning, Computer Vision, Machine Learning, Computer System Organization (C), Algorithms, Data Structures, Functional Programming (OCaml), Discrete Math, Linear Algebra
  \item \textbf{Honors:} BURE Fellow, CIDA Grant Recipient, AIMI Junior Affiliate, 4-time AIME Qualifier, USACO Silver Division
\end{itemize}

\section{Technical Skills}
\begin{itemize}
  \item \textbf{Languages:} Python, Java, C/C++, TypeScript/JavaScript, OCaml
  \item \textbf{ML:} PyTorch, Hugging Face Transformers, PEFT, scikit-learn, NumPy, pandas, OpenCV, YOLO
  \item \textbf{Tools:} Git, GitHub Actions, SLURM/HPC, CUDA, Playwright, React, Node.js, Claude API
\end{itemize}

\section{Research Experience}
\entry{Undergraduate Researcher, Sun Lab at Cornell University}{May 2025 -- Present}{PI: Jennifer Sun. AI for science: scientific discovery agents and computer vision for animal behavior}{}
\begin{itemize}
  \item \textbf{Statistical validation for automated hypothesis discovery (AIMI Junior Affiliate), Jun 2026 -- Present.} Implemented the randomization tests (p- and e-values for 11 claim types) and the false discovery rate reporting rules behind an independent check on what AI discovery agents report
  \item Ran five agent systems, 500+ runs in 13 test settings with known ground truth: the agents' own reports were 18 to 66\% false; the check brought that below 1\% for every agent; stress tests showed it can fail if held-out results leak
  \item \textbf{Multi-lab video behavior recognition (MABe), Feb 2026 -- Present.} Beat the frozen-backbone baseline on 9 of 15 lab datasets (13 of 15 on at least one split) by LoRA fine-tuning a video foundation model with a self-authored kinematic auxiliary loss; on one lab, video alone beat the pose-tracking ceiling
  \item Wrote and ran the 8-stage SLURM chain that took all 15 labs from raw video to evaluated 5-fold ensembles, staging up to 1.9 TB of features per lab on node-local NVMe; ran on Blackwell GPUs
  \item Moved the VideoPRISM backbone from JAX to PyTorch with parity checked, which enabled LoRA fine-tuning; wrote the PyTorch training stage, video loader, and 21 unit tests
  \item \textbf{Calf behavior monitoring (BURE Fellow, CIDA grant), May 2025 -- Present.} Mean posture-classification F1 rose from 0.928 with 300 training images to 0.957 with 1,000 (24 runs per setting, DINOv2 vs.\ DINOv3 ensembles); fine-tuned YOLO on calf data I annotated in CVAT. Poster presented at ADSA 2026; manuscript in preparation
  \item Benchmarked five vision-language models (Qwen, Gemma 3, Gemini 2.5 Pro, PaliGemma, NVLM) for ear-tag and pose recognition; built a Jupyter workflow so veterinary researchers can run experiments themselves
\end{itemize}

\section{Publications}
\begin{itemize}
  \item \textbf{R. M. Ye}, M. Madureira Ferreira, J. J. Sun, F. A. Leal Yepes. ``Automating pose detection in calves: a computer vision framework for non-invasive behavioral monitoring.'' Poster, ADSA Annual Meeting, 2026
  \item M. Madureira Ferreira, K. M. Dadallage, T. Ward, H. McCray, \textbf{R. M. Ye}, J. J. Sun, F. A. Leal Yepes. ``A smarter lung ultrasonography solution for lung consolidation classification in dairy calves using a deep learning approach.'' \textit{Journal of Dairy Science}, 2026. Supporting author
\end{itemize}

\section{Projects}
\entry{DoRA Re-Implementation and Evaluation}{Spring 2026}{CS 4782 Deep Learning final project, 4-person team. PyTorch, Hugging Face PEFT}{\href{https://github.com/leonjiao5/cs4782-project}{GitHub}}
\begin{itemize}
  \item Team re-implementation of DoRA, tested against LoRA on Qwen2.5-3B (AMC12, MATH); ran all HPC experiments
  \item Built the evaluation harness (logit-probe and two-pass scoring, separating what a model knows from what it generates), the data pipeline with a leakage filter, and the figures; diagnosed a critical bug in the DoRA layer
\end{itemize}

\entry{Worship Rig}{Sep 2026}{Personal project. Electron, JavaScript, Web Audio}{\href{https://ryanmye.github.io/worship-rig/}{Site} $\cdot$ \href{https://github.com/ryanmye/worship-rig}{GitHub}}
\begin{itemize}
  \item Built and shipped a Mac app that turns a MIDI keyboard into a keys rig for worship bands: 43 instruments, a key drone, setlists, per-sound EQ, and WAV recording, on a Web Audio engine that also renders offline for tests
  \item Wrote 1,000+ automated tests (Playwright), a GarageBand sample converter, and CI builds for Apple Silicon and Intel
\end{itemize}

\entryshort{\href{https://github.com/matthewj5/sisyphus}{Sisyphus: AI Productivity App} (Claude Builders Club Hackathon, 4-person team)}{Nov 2025}
\begin{itemize}
  \item Checks task photos with Claude; one of 5 winners among 38 projects (Fun Award); worked on frontend and AI validation
\end{itemize}

\section{Leadership}
\noindent\textbf{Basic Coding}, Co-Founder \& President. Free coding classes for 300+ students in 8 states \hfill Aug 2020 -- Jun 2024

\end{document}
